\section{A central dihedral obstruction in dimension 120}
\label{sec:rank11-d120}

\paragraph{Theorem.}
The based rings \texttt{R11\_1\_40} and \texttt{R11\_1\_82} admit no
categorification by a fusion category over $\mathbb C$.

\paragraph{Exact hypotheses.}
Both rings are commutative of rank eleven, with dimensions
$1^2,2,3^2,4^6$. Their nontrivial invertible $g$ has order two;
its multiplication permutation exchanges the two invertibles, exchanges
the two dimension-three simples, and fixes the other seven simples.
The basis elements of dimension less than four form a closed subring
of dimension twenty-four. No identification of this subring with a
particular category is assumed: in the second input its two dimension-three
simples are dual to each other, unlike the representation ring of $S_4$.

For a scalar character $\chi$, put
$f_\chi=\sum_i\chi(b_i)\chi(b_i^*)$ and $\beta_\chi=120/f_\chi$.
The eleven values $\beta_\chi$ are
\[
 1,2,2,8,8,8,8,8,15,30,30.
\]
Both characters with $\beta_\chi=2$ satisfy $\chi(g)=1$.
Finally the induction Gram coefficient $B_{gg}$ is eleven.
The accompanying exact verifier checks these statements directly from
the source tensor and a complete character table, including every
character multiplication identity. No restriction-frame search is used.

\paragraph{A forced central dihedral category.}
Suppose that $\mathcal C$ is a categorification. Its integrality supplies
the canonical positive spherical structure by
\cite[Proposition~8.24]{ENO}. Let $F:Z(\mathcal C)\to\mathcal C$ be
forgetful, with right adjoint $I$. By \cite[Theorems~2.5 and~2.13]{OstrikPivotal},
the simple constituents $A_\chi$ of $I(\mathbf1)$ occur once, have
$\dim(A_\chi)=\beta_\chi$, satisfy $[F(A_\chi):\mathbf1]=1$, and have
trivial twist. The character identification preceding Theorem~2.13 gives
\[
 S_{A_\chi,Y}=\beta_\chi\chi(F(Y))
 \qquad(Y\text{ simple in }Z(\mathcal C)),
\]
where $S$ is unnormalized. For the two dimension-two constituents $A,B$,
positivity forces $F(A)=F(B)=\mathbf1\oplus g$, whence
\[
 S_{A,A}=S_{A,B}=S_{B,B}=4.
\]
These equalities force trivial monodromy: the balancing formula is a
positive weighted sum of roots of unity, and equality with the sum of
the positive weights forces every eigenvalue to be one. Therefore the
subcategory $\mathcal T=\langle A,B\rangle$ is symmetric and has trivial
twist. It is Tannakian; see also \cite{DGNO}.

All objects of $\mathcal T$ forget to sums of $\mathbf1$ and $g$.
If a simple central $Y$ with this property contains $\mathbf1$ after
forgetting, it lies in $I(\mathbf1)$, with unit multiplicity one. Write
$F(Y)=\mathbf1\oplus m g$. The induction Gram equality gives
$m^2\le B_{gg}=11$, so $\dim(Y)\le4$. The displayed $\beta$ list
therefore leaves only $Y=\mathbf1,A,B$. In particular, the only simple
central object supported solely on $\mathbf1$ is $\mathbf1$ itself.

If instead $F(Y)=m g$, then $F(Y^*Y)=m^2\mathbf1$; all constituents of
$Y^*Y$ must be the central unit. Simplicity of $Y$ gives $m^2=1$.
Such a $Y$ is invertible and forgets to $g$. There is at most one such
lift, since the ratio of two lifts forgets to the unit; call it $z$ if
it exists. It must occur in $\mathcal T$: the dimension-four object
$A^*A$ contains the unit exactly once, and its remaining dimension three
cannot be filled using only the dimension-two simples $A,B$.
Thus $\mathcal T$ has exactly $\mathbf1,z,A,B$, of dimensions $1,1,2,2$.
It is consequently $\operatorname{Rep}(D_{10})$, where $D_{10}$ has order
ten. Its unique nontrivial invertible $z$ is a central boson with $F(z)=g$.

\paragraph{De-equivariantization and a faithful grading.}
The subcategory $\mathcal E=\langle z\rangle\simeq\operatorname{Rep}(C_2)$
embeds in $\mathcal C$ under $F$. By
\cite[Proposition~2.10]{ENOSolvable}, de-equivariantization gives a
category $\mathcal D$ of dimension sixty with
$\mathcal C\simeq\mathcal D^{C_2}$.
A free orbit $\{X,gX\}$ gives one simple of dimension $\dim(X)$;
a fixed $X$ gives two simples of dimension $\dim(X)/2$.
Here the fixed stabilizer algebra has two one-dimensional representations,
since $H^2(C_2,\mathbb C^\times)=0$, and its sign automorphism interchanges
them. The resulting dimension list is
\[
 \mathcal D:\quad 1^3,2^{12},3,
 \qquad \operatorname{rank}(\mathcal D)=16.
\]

Since $\mathcal T$ is symmetric, it lies in $\mathcal E'$.
The central equivalence $\mathcal E'_{C_2}\simeq Z(\mathcal D)$ sends
$\mathcal T_{C_2}$ to $\operatorname{Rep}(C_5)$.
Every object of this subcategory forgets to a multiple of the unit:
$\mathbf1,g$ become the unit under de-equivariantization, and the images
of $A,B$ become two copies of it. Their four nontrivial $C_5$-character
summands must each forget to the unit by positivity.
Proposition~2.9(i) of \cite{ENOSolvable} now supplies a faithful $C_5$
grading of $\mathcal D$, whose components all have dimension twelve.

The closed dimension-twenty-four subcategory above de-equivariantizes
to a subcategory $\mathcal K\subset\mathcal D$ of dimensions $1,1,1,3$.
Its invertibles form $C_3=\langle a\rangle$, and its unique other simple
$X$ is self-dual and satisfies
\[
 aX=Xa=X,\qquad X^2=\mathbf1+a+a^2+2X.
\]
Thus $\mathcal K$ has the fusion ring of $\operatorname{Rep}(A_4)$.
It equals its adjoint category, so every grading restricts trivially to
$\mathcal K$. Equality of dimensions gives $\mathcal D_e=\mathcal K$.

\paragraph{The Morita class of every possible neutral category.}
By \cite[Theorem~1.6]{ENOSolvable}, $\mathcal K$ is solvable. A solvable
tower consists of extensions and equivariantizations by cyclic groups of
prime order \cite[Definition~1.2 and Proposition~4.4]{ENOSolvable}.
Its final step cannot be an extension because $\mathcal K$ has trivial
universal grading. Hence $\mathcal K=\mathcal L^{C_p}$ for a prime $p$.
The embedded $\operatorname{Rep}(C_p)$ forces $p=3$, and
$\operatorname{FPdim}(\mathcal L)=4$.

The simple equivariant dimensions are $[C_3:H]\dim(\pi)\dim(Y)$.
All projective irreducible representations of the cyclic stabilizers $H$
are one-dimensional. Since the only dimensions in $\mathcal K$ are one
and three, a fixed $Y$ has dimension one or three, the latter impossible
in a category of dimension four. A nonfixed $Y$ has an orbit of length
three and likewise has dimension one. Thus
$\mathcal L=\operatorname{Vec}_H^\eta$ for a group $H$ of order four.
The action on $H$ is nontrivial, since otherwise equivariantization would
have twelve invertibles. Thus $H=V_4$, with $C_3$ cycling its nonidentity
elements.

The pointed crossed product $\mathcal L\rtimes C_3$ is Morita equivalent
to $\mathcal K$ by \cite[Proposition~2.2]{ENOSolvable}. Its group of
invertibles is the extension of $V_4$ by $C_3$ with this nontrivial action;
the coprime extension splits, giving $A_4$. Consequently
\[
 \mathcal K\text{ is Morita equivalent to }\operatorname{Vec}_{A_4}^\Omega
\]
for some three-cocycle $\Omega$, with no triviality assumption on $\Omega$.

\paragraph{No order five in the twisted Brauer--Picard group.}
Let $\mathcal P=\operatorname{Vec}_{A_4}^\Omega$.
The twisted classification \cite[Theorem~1.2, or Theorem~4.12]{NaiduNikshych}
parametrizes Lagrangian subcategories of $Z(\mathcal P)$ by normal abelian
subgroups and invariant alternating $\Omega$-bicharacters.
The only possible subgroups are $1,V_4$. The first contributes one choice;
for $V_4$, the ratio of two admissible twisted bicharacters is an ordinary
alternating bicharacter, of which there are exactly two. Hence there are
at most three Lagrangian subcategories.

The stabilizer of the canonical $\operatorname{Rep}(A_4)$ preserves its
regular algebra $\operatorname{Fun}(A_4)$: a connected etale algebra of
dimension twelve in $\operatorname{Rep}(A_4)$ corresponds uniquely to the
free transitive $A_4$-set. Conversely the regular algebra generates that
subcategory. Theorem~3.3 and Remark~3.4 of \cite{MarshallNikshych} identify
this stabilizer with the induction image $I(A_4,\Omega)$.
The elementary bound from tensor autoequivalences already suffices:
their kernel over the underlying automorphism group is
$H^2(A_4,\mathbb C^\times)=C_2$ and their image is contained in
$\operatorname{Aut}(A_4)=S_4$, so the stabilizer order divides $48$.

More precisely, Proposition~4.5 of \cite{MarshallNikshych}, valid for
arbitrary $\Omega$, gives
\[
 Z(A_4)\longrightarrow H^2(A_4,\mathbb C^\times)
 \longrightarrow I(A_4,\Omega)
 \longrightarrow\operatorname{OutStab}(\Omega)\longrightarrow1.
\]
Since $Z(A_4)=1$ and $\operatorname{Out}(A_4)=C_2$, the stabilizer order
is two or four. These group-theoretic data are recorded in
\cite[Section~8.3]{NikshychRiepel}; its untwisted Brauer--Picard computation
is not being applied to a twisted category.
Orbit--stabilizer, with orbit size at most three, shows that
$\operatorname{BrPic}(\mathcal P)\simeq
\operatorname{Aut}_{\mathrm{br}}(Z(\mathcal P))$ has no element of order
five. Morita invariance gives the same statement for $\mathcal K$.

\paragraph{Contradiction.}
A faithful graded extension determines a homomorphism
\[
 C_5\longrightarrow\operatorname{BrPic}(\mathcal K),
 \qquad h\longmapsto[\mathcal D_h].
\]
It is trivial, so every component is equivalent to the regular
$\mathcal K$-bimodule category and has rank four. Thus
$\operatorname{rank}(\mathcal D)=20$, contrary to the forced rank sixteen.
This proves the theorem for both inputs.\hfill$\square$
